\begin{figure}[tb]
\centering
\begin{tikzpicture}[>=Latex,font=\small]
  \begin{scope}[scale=.48]
    \draw[->,black!55] (-.45,0) -- (13.1,0)
      node[below=2pt,text=black] {\(e/p\)};
    \draw[->,black!55] (0,-.35) -- (0,5.0)
      node[above,text=black] {\(d/p\)};
    \fill[blue!8] (0,0) -- (.75,1) -- (12,4) -- cycle;
    \draw[blue!60!black,thick] (0,0) -- (.75,1) -- (12,4) -- cycle;
    \fill (0,0) circle (2.4pt) node[below left] {\(v_1\)};
    \fill (.75,1) circle (2.4pt) node[above left] {\(v_2\)};
    \fill (12,4) circle (2.4pt) node[above right] {\(v_3\)};
    \fill[orange!75!black] (4.25,1.666667) circle (2.4pt);
    \draw[orange!85!black,thick] (4.25,1.666667) circle (.115);
    \node[above left] at (4.0,2.25) {\(g\)};
    \draw[black!55] (3.93,2.12) -- (4.18,1.80);
    \draw[black!55,dashed] (3.75,1.15) rectangle (4.75,2.2);
  \end{scope}
  \draw[black!55,->] (2.28,1.05) to[bend left=14] (7.48,2.3);
  \begin{scope}[shift={(8.9,1.65)}]
    \fill[orange!8] (0,0) circle (1.35);
    \draw[orange!85!black,thick] (0,0) circle (1.35);
    \draw[black!55,dashed] (-1.7,0) -- (1.7,0);
    \draw[very thick,blue!65!black] (-.93,0) -- (.93,0);
    \fill[blue!65!black] (-.93,0) circle (2pt)
      node[below=5pt] {\((s_1,\rho)\)};
    \fill[blue!65!black] (.93,0) circle (2pt)
      node[below=5pt] {\((s_2,\rho)\)};
    \fill (0,0) circle (1.6pt) node[above=4pt] {\(g\)};
    \node at (0,1.67) {\(\overline B(g,\gamma)\)};
    \node at (0,-1.72) {\(\rho=5p/3,\quad s_1<s_2\)};
  \end{scope}
\end{tikzpicture}
\caption{Pressure-support coordinates \((e,d)\), with the left panel
drawn in coordinates \((e/p,d/p)\).
The triangle of the initial limiting slopes \(v_i\) contains a ball
about its barycenter \(g\);
the ball persists under the controlled perturbations of its vertices.
The schematic inset shows two points on the central horizontal slice
at density \(5p/3\).  Once the ball lies in the subdifferential of a
pressure at one parameter, each nearby horizontal slice supplies two
supports with distinct inverse-temperature slopes at that same parameter.
Lemma~\ref{m:canonical-corner} converts this geometry into the canonical
derivative jump throughout a density interval.
}
\label{fig:support-geometry}
\end{figure}
